\documentclass[11pt]{article}
\usepackage[margin=1in]{geometry}

% Core packages
\usepackage{amsmath,amssymb}
\usepackage{tikz-cd}
\usepackage{multicol}

% Paragraphs
\setlength{\parindent}{0pt}
\setlength{\parskip}{1\baselineskip}

\title{DISTANCE v0.1\\[0.5em]
\large A Philosophical Framework for the Direction of the Universe}
\author{Kangbin Yim\\Soonchunhyang University}
\date{2024}

\begin{document}
\maketitle

For a long time, humanity has understood the universe as a collection of objects placed within space. Stars float in space, planets move between them, and life emerges and disappears at specific locations. We perceive time as something that flows, space as an empty background, and reality as something explainable through distance, speed, position, and mass. These assumptions have become so familiar that we rarely question them.

Newton interpreted time and space as absolute frameworks, while Einstein reinterpreted them as relative structures. Yet even after modern physics transformed our understanding of reality, we still continued to treat time and space themselves as fundamental substances of the universe. We assumed that events occur ``within time'' and that objects move ``through space.''

But at some point, I began asking a different question.

What if time and space are not the deepest layer of reality?
What if they are not the true hardware of the universe, but merely interfaces constructed by human cognition in order to interpret something more fundamental?

When we observe the second hand of a clock moving, we say that time has passed. Yet what actually occurred was energy transfer. A spring unwound, a chemical reaction generated current, a circuit operated, and the position of the hand changed. What we call ``time'' may simply be an interpretation applied after a change has already occurred.

Space may be similar. When we say that an object moved from one place to another, what may actually have occurred is not movement through an independent spatial container, but rather a reconfiguration of energy states and relationships. Human beings may merely translate those changes into the language of distance and position.

Gradually, I began to see the universe not as a collection of independent entities, but as a continuous transformation of energetic relationships.

Whenever two entities possess different energy states, the possibility of change emerges between them. Energy does not remain perfectly uniform; it tends to flow toward lower states. I call this possibility and directionality of relational transformation ``momentum.''

Here, momentum does not simply mean physical momentum in the conventional sense. It refers more fundamentally to the potential for relational change generated by differences in energy states between entities. All transformations in the universe begin from this momentum.

And it was through this perspective that I arrived at the concept of ``Distance.''

Distance, in this context, is not merely spatial separation or visual discontinuity. Distance is an asymmetrical relational state formed through differences in energy conditions. It is the condition in which directionality and intensity of energy transfer emerge between entities of differing energetic states.

When momentum is large, Distance becomes small.
When momentum is weak, Distance becomes large.

Seen this way, the universe is no longer a collection of objects suspended in space. It becomes a field of Distances---a vast dynamic structure in which energetic differences generate momentum, momentum generates relationships, and relationships continuously reorganize themselves.

Time becomes the way human beings perceive this flow.
Space becomes an interface through which human beings translate relational structures into understandable coordinates.

Existence itself must then be reconsidered.

We usually perceive humans, stars, organisms, nations, and civilizations as independently existing objects. But perhaps existence is not an independent substance at all. Perhaps what we call existence is merely an image temporarily drawn by a particular energetic condition.

An entity may not be a fixed object, but rather a transient flow-state maintained within ongoing transformations of Distance. Each existence continuously alters the momentum of others, rearranges relationships, and contributes to the entropic direction of the universe. In this sense, entities function as local engines of Distance transformation.

Life is no exception.

For a long time, life has been interpreted as something resisting entropy because living systems maintain internal order, reproduce themselves, and stabilize structure. Yet from another perspective, life may not resist the universal flow at all. Life may instead be a self-organizing structure that explores, connects, amplifies, and redistributes Distance.

Plants utilize energetic differences between stars and planets. Animals move along chemical gradients. Microorganisms decompose collapsed structures and generate new flows. Life is not an exception standing outside the universe, but a localized vortex formed within the flow of Distance itself.

Humanity is not the center of this process. Human civilization is simply one large-scale composite structure generated by the flow of Distance. Cities, industries, networks, and information systems all function by connecting and redistributing enormous energetic asymmetries.

Civilization appears to create order, but at the same time it accelerates entropy and Distance transformation on a much larger scale.

Entropy itself must therefore be redefined.

Entropy is not merely the increase of disorder. It is the directional increase of Distance between entities. Momentum weakens, relational intensity decreases, and the possibility of energetic transfer gradually fades.

The ultimate fate of the universe may not simply be thermal death in the traditional sense, but a condition in which Distance between all entities becomes effectively infinite---a state where no further relational transformation is possible. Only minimal energy-bearing subjects remain, separated by infinite Distance, with no momentum left capable of generating change.

Even the Big Bang may need reinterpretation.

Perhaps the Big Bang was not simply an expansion of physical space, but the beginning of increasing Distance emerging from an extreme state of near-zero Distance and maximal momentum. The universe may not be expanding spatially so much as relationally drifting apart.

Yet this work is not an attempt to establish another absolute cosmology. Rather, it is a philosophical suspicion directed toward humanity's confidence that it has already understood reality correctly.

We may have mistaken interfaces for reality itself.

DISTANCE is therefore less an answer than a different way of questioning.

Does the universe truly exist inside space?
Are entities truly independent substances?
Does time actually flow?
Is life really an exception to entropy?
And why have human beings become so certain about interpreting reality in the way we do?

This work begins from those questions.

\end{document}